% Çıkarım hedefi, ekvator düzleminde (z = 0) kesilmiş.
% Yalnız çizgi işi: alt yarı için tel kafes, iç yapı için dolu kesit yüzü.
% Her sayı `../preprint/veri/`den okunur — bkz. scripts/onkayit_veri.py.
% `docs/rapor2204` içinden derle:  lualatex sekil/hedef.tex
\documentclass[border=5pt]{standalone}
\input{sekil/ortak.tex}
\usepackage{tikz-3dplot}

\input{\VERI/carpma.tex}
\pgfplotstableread{\VERI/blok_kesit_z0.dat}\blokkesit

% `etiket` biçemi ortak.tex'te \footnotesize sabitliyor ve resim düzeyindeki
% font ayarını eziyor; bu şekle özel olarak büyütülür (yukarıdaki gerekçe).
\tikzset{etiket/.style={inner sep=1pt, font=\Large}}

\tdplotsetmaincoords{64}{28}
% 88,5 m -> 7 cm, yani 177 m'lik cisim sayfada ~14 cm yer kaplar.
\pgfmathsetmacro{\olcek}{7.0/88.5}

\begin{document}
% Yazı tipi ön baskıdakinden büyük (\footnotesize -> \Large): bu şekil
% doğal olarak 244 mm genişliğinde ve A4 tek kolonda 0,64 kat KÜÇÜLTÜLEREK
% basılıyor, o oranda \footnotesize gövde metninin yanında ~5 pt'ye düşüyordu.
\begin{tikzpicture}[font=\Large, >={Stealth[length=2.2mm,width=1.6mm]}]
\begin{scope}[tdplot_main_coords, scale=\olcek]

\def\A{88.5} \def\B{87.0} \def\C{58.0}

% ---------- 1. alt yarının tel kafesi (z <= 0) --------------------------
\foreach \v in {15,32,50,70}{%
  \draw[draw=tel, line width=2.0pt, opacity=0.38]
    plot[variable=\t, domain=0:360, samples=120, smooth cycle]
    ({\A*cos(\v)*cos(\t)}, {\B*cos(\v)*sin(\t)}, {-\C*sin(\v)});
}
\foreach \u in {0,30,...,330}{%
  \draw[draw=tel, line width=2.0pt, opacity=0.38]
    plot[variable=\t, domain=0:90, samples=48, smooth]
    ({\A*cos(\t)*cos(\u)}, {\B*cos(\t)*sin(\u)}, {-\C*sin(\t)});
}

% ---------- 2. kesit yüzü, opak: arkasındaki kafesi gizler -------------
\draw[fill=yuzey, draw=koyu, line width=3.0pt]
  plot[variable=\t, domain=0:360, samples=180, smooth cycle]
  ({\A*cos(\t)}, {\B*sin(\t)}, 0);

% ---------- 3. kesit düzleminin kestiği bloklar ------------------------
% Gölgeli, yassı değil: bunlar kürelerden alınmış dilimler ve düz dolgu
% onları gözleme gibi gösteriyordu.
\pgfplotstablegetrowsof{\blokkesit}
\pgfmathtruncatemacro{\nb}{\pgfplotsretval-1}
\foreach \i in {0,...,\nb}{%
  \pgfplotstablegetelem{\i}{x}\of\blokkesit \let\bx\pgfplotsretval
  \pgfplotstablegetelem{\i}{y}\of\blokkesit \let\by\pgfplotsretval
  \pgfplotstablegetelem{\i}{r}\of\blokkesit \let\br\pgfplotsretval
  \shadedraw[inner color=blokic, outer color=blokdis,
             draw=koyu!65, line width=2.6pt]
    plot[variable=\t, domain=0:360, samples=90, smooth cycle]
    ({\bx+\br*cos(\t)}, {\by+\br*sin(\t)}, 0);
}

% ---------- 4. incelme kabukları, aynı düzlemle kesilmiş ---------------
\begin{scope}
  \clip plot[variable=\t, domain=0:360, samples=180, smooth cycle]
    ({\A*cos(\t)}, {\B*sin(\t)}, 0);
  \foreach \r in {48,24,12,6,3}{%
    \draw[draw=kabuk, line width=3.4pt, opacity=0.9]
      plot[variable=\t, domain=0:360, samples=140, smooth cycle]
      ({\carpmaX+\r*cos(\t)}, {\r*sin(\t)}, 0);
  }
\end{scope}

% ---------- 5. çarpma --------------------------------------------------
\coordinate (I) at (\carpmaX, 0, 0);
\coordinate (D) at ({\carpmaX-62*\yonX}, 0, {-62*\yonZ});
\coordinate (N) at ({\carpmaX+34}, 0, 0);
\draw[densely dashed, koyu!55, line width=2.6pt] (I) -- (N);
\draw[vurgu, line width=6pt, ->, >={Stealth[length=15pt,width=11pt]}]
      (D) -- (I);
\fill[vurgu] (I) circle (3.4pt);

\coordinate (ucA) at ({\carpmaX+48}, 0, 0);
\coordinate (ucB) at ({\carpmaX-48}, 0, 0);
\end{scope}

% ---------- 6. etiketler DIŞARIDA, kılavuz çizgilerin ucunda -----------
% bloklar — sol üstteki bloğa kılavuz
\node[etiket, anchor=south east, align=right, text=blokdis!75!black]
     (lblok) at ($(I)+(-12.6,3.1)$)
     {bloklar --- gözeneklilik $\alpha_b$, hacim kesri $f$\\
      yarıçaplar $r\in[14,56]\ \mathrm{m}$,
      $\mathrm{d}N/\mathrm{d}r\propto r^{-3}$};
\draw[kilavuz=blokdis] (lblok.south east) -- ($(I)+(-10.9,1.55)$);

% matris — soldaki boş matrise kılavuz
\node[etiket, anchor=east, text=koyu!70] (lmat) at ($(I)+(-13.2,0.1)$)
     {matris --- bağdaşım dayanımı $Y_0$};
\draw[kilavuz=koyu!60] (lmat.east) -- ($(I)+(-11.6,-0.3)$);

% incelme kabukları — sağdan, yörüngenin üstünden
\node[etiket, anchor=west, align=left, text=kabuk] (lkab)
     at ($(I)+(1.5,1.5)$)
     {iç içe incelme kabukları\\$r=3{,}\,6{,}\,12{,}\,24{,}\,48\ \mathrm{m}$};
\draw[kilavuz=kabuk] (lkab.west) -- ($(I)+(-0.75,0.55)$);

% DART — yörüngenin altında, oktan uzakta
\node[etiket, anchor=north west, align=left, text=vurgu]
     at ($(I)+(3.05,-2.75)$)
     {\textbf{DART}\\[1pt]
      $6144{,}9\ \mathrm{m\,s^{-1}}$, $579{,}4\ \mathrm{kg}$\\
      yüzey normalinden $\aciDeg^{\circ}$};

% Şeklin içine gömülü altyazı yok: rapor içinde altyazıyı LaTeX verir,
% ikisi birden olunca altyazı iki kez yazılmış gibi görünüyordu.
\end{tikzpicture}
\end{document}
